% Chapter 2: Proto-cohesion Diagnostic Vector
% Part I: Foundations

\chapter{Proto-cohesion 진단 벡터:\\Bind, Sep, Inside, Persist}
\label{ch:diagnostic-vector}

\begin{quote}
\textit{``형성의 질을 묻는 것은 하나의 숫자가 아니라 네 개의 독립적 차원에서 묻는 것이다.''}
\end{quote}

\section{Soft Cohesion Field의 소개}
\label{sec:soft-field-intro}

Chapter~\ref{ch:motivation}에서 우리는 SCC 이론이 부드러운 응집 장(soft cohesion field) $u_t : X_t \to [0,1]$을 원초적 존재자로 삼는다는 것을 보았다. 이 장에서는 이 응집 장이 주어졌을 때, 그것이 얼마나 ``좋은 형성(good formation)''인지를 평가하는 도구를 소개한다.

\subsection{$[0,1]^X$ 필드의 의미}

관계적 지지 공간 $X_t = \{x_1, x_2, \ldots, x_n\}$이 $n$개의 지점으로 이루어져 있을 때, 응집 장은 $n$차원 벡터이다:
\begin{equation}
    u_t = (u_t(x_1), u_t(x_2), \ldots, u_t(x_n)) \in [0,1]^n
    \label{eq:field-as-vector}
\end{equation}

이 벡터는 각 지점의 응집적 참여 강도를 부호화한다. 그러나 모든 $[0,1]^n$ 벡터가 의미 있는 형성을 표현하는 것은 아니다. 임의의 벡터---예를 들어 무작위로 생성된 값들---는 내부적 일관성이 없고, 외부와 구별되지 않으며, 시간적으로 지속하지 않는다.

``좋은'' 응집 장, 즉 진정한 응집적 형성을 표현하는 장은 네 가지 독립적 구조적 요구를 만족해야 한다:
\begin{enumerate}
    \item \textbf{결합(Binding)}: 장이 관계적으로 자기-지지적인가?
    \item \textbf{분리(Separation)}: 장이 외부와 구조적으로 구분되는가?
    \item \textbf{내부성(Inside)}: 장이 일관된 형태학적 구조를 가지는가?
    \item \textbf{지속성(Persistence)}: 장이 시간적으로 계승되는가?
\end{enumerate}

이 네 가지 질문에 대한 연속적 답변이 바로 \textbf{proto-cohesion 진단 벡터}이다.


\subsection{부피 제약}

응집 장에는 하나의 전역적 제약이 부과된다. 모든 지점의 응집 값의 합이 일정해야 한다:
\begin{equation}
    \sum_{x \in X_t} u_t(x) = m
    \label{eq:volume-constraint}
\end{equation}
여기서 $m > 0$은 총 응집 부피이다. 부피 분율 $c = m/n$은 필드의 평균값이다.

이 제약은 기술적 필요성과 존재론적 자연스러움을 모두 가진다. 기술적으로, 이 제약이 없으면 모든 에너지 항이 0인 자명한 해 $u \equiv 0$이 전역 최소화자가 된다. 존재론적으로, 이 제약은 유한한 응집 용량(finite cohesive capacity)에 대한 약속이다---인지 시스템이 유한한 자원을 가진다는 자연스러운 가정.

제약 하에서 장이 사는 공간을 \textbf{부피 제약 단체(volume-constrained simplex)}라 한다:
\begin{equation}
    \Sigma_m = \Big\{u \in [0,1]^n\ \Big|\ \sum_{i=1}^n u_i = m\Big\}
    \label{eq:simplex}
\end{equation}


\section{진단 벡터 개관}
\label{sec:diagnostic-overview}

Proto-cohesion 진단 벡터는 네 개의 $[0,1]$-값 요소로 구성된다:
\begin{equation}
    \boxed{\mathbf{d} = (\mathsf{Bind},\ \mathsf{Sep},\ \mathsf{Inside},\ \mathsf{Persist}) \in [0,1]^4}
    \label{eq:diagnostic-vector}
\end{equation}

각 요소는 형성의 한 구조적 측면을 연속적으로 평가한다. 네 요소 모두 1에 가까운 형성은 ``완전한 객체''에 가깝다. 일부만 높은 형성은 부분적으로 응집된 전-객체(pre-object)이다. 진단 벡터는 불리언(Boolean) 판정이 아니라 연속적 평가이므로, ``경계 사례''를 자연스럽게 다룬다.

\begin{table}[t]
\centering
\caption{Proto-cohesion 진단 벡터의 네 요소 개관}
\label{tab:diagnostic-overview}
\begin{tabular}{lllc}
\hline
\textbf{요소} & \textbf{평가 대상} & \textbf{직관} & \textbf{에너지 대응} \\
\hline
$\mathsf{Bind}$ & 관계적 자기-지지 & ``장이 스스로를 유지하는가?'' & $\mathcal{E}_{\mathrm{cl}}$ \\
$\mathsf{Sep}$  & 외부와의 구조적 대비 & ``내부와 외부가 구분되는가?'' & $\mathcal{E}_{\mathrm{sep}}$ \\
$\mathsf{Inside}$ & 형태학적 분절 & ``핵심부-경계-외부가 명확한가?'' & $\mathcal{E}_{\mathrm{bd}}$ \\
$\mathsf{Persist}$ & 시간적 구조 계승 & ``이전 형성을 계승하는가?'' & $\mathcal{E}_{\mathrm{tr}}$ \\
\hline
\end{tabular}
\end{table}


이제 각 요소를 하나씩 살펴보자. 각 요소에 대해 먼저 \emph{직관}을 설명하고, 그 다음 \emph{수학적 정의}를 제시하며, 마지막으로 구체적 \emph{예제}에서 수치를 계산한다.


\section{Bind: 관계적 자기-지지}
\label{sec:bind}

\subsection{직관}

``이 형성은 자기 자신을 유지할 수 있는가?''

Bind(결합)는 응집 장이 닫힘 연산자 $\mathrm{Cl}_t$에 의한 관계적 완성과 얼마나 일치하는지를 측정한다. 직관적으로, 닫힘 연산은 각 지점의 응집값을 이웃들로부터 받는 관계적 지지에 맞게 갱신한다. 만약 장이 이미 자기-지지적이라면---갱신해도 변하지 않는다면---그 장은 완벽하게 결합(bound)되어 있다.

비유하자면: 건물의 각 기둥이 인접 기둥들과 상호 지지하는 상태를 생각하자. 각 기둥이 정확히 이웃들로부터 받는 지지만큼의 하중을 견디고 있다면, 그 구조는 자기-지지적이다. Bind는 이러한 자기-지지의 정도를 측정한다.

% [Figure placeholder: Bind visualization]
\begin{figure}[t]
\centering
% \includegraphics[width=0.8\textwidth]{figures/ch02_bind_visualization.pdf}
\fbox{\parbox{0.8\textwidth}{\centering\vspace{3.5cm}
\textbf{[그림 2.1]} Bind 시각화.\\
좌: 원본 장 $u$ (Gaussian bump). 중앙: 닫힘 적용 후 $\mathrm{Cl}(u)$.\\
우: 잔차 $|u - \mathrm{Cl}(u)|$ (경계 부근에서 가장 큼).\\
Bind = $1 - \|u - \mathrm{Cl}(u)\|_2 / \sqrt{n}$.
\vspace{0.5cm}}}
\caption{Gaussian bump 형성에 대한 Bind 계산. 잔차(우측)는 경계 대역에서 가장 크다. 이는 경계 지점이 내부 닫힘력과 외부 이완 사이의 구조적 긴장 상태에 있기 때문이다.}
\label{fig:bind-visualization}
\end{figure}


\subsection{수학적 정의}

\begin{equation}
    \boxed{\mathsf{Bind}_t(u_t) = 1 - \frac{\|u_t - \mathrm{Cl}_t(u_t)\|_2}{\sqrt{n}}}
    \label{eq:bind-def}
\end{equation}

여기서:
\begin{itemize}
    \item $\mathrm{Cl}_t(u_t)$는 닫힘 연산자를 $u_t$에 적용한 결과
    \item $\|\cdot\|_2$는 $\ell^2$ (유클리드) 노름
    \item $n = |X_t|$는 지지 공간의 크기
    \item $\sqrt{n}$으로 나누어 크기 독립성(scale-independence)을 보장
\end{itemize}

\textbf{해석}:
\begin{itemize}
    \item $\mathsf{Bind} = 1$: 장이 닫힘의 고정점 ($u = \mathrm{Cl}(u)$)---완벽한 자기-지지
    \item $\mathsf{Bind} \approx 0$: 장과 닫힘 결과가 크게 다름---자기-지지 실패
\end{itemize}

\textbf{왜 $\ell^2$ 노름인가?} $\ell^\infty$ 노름은 경계 지점의 본질적 잔차(inherent residual)로 인해 부적합하다. 이중 우물 포텐셜과 닫힘 사이의 구조적 긴장 때문에, 경계 지점에서는 $|u - \mathrm{Cl}(u)| \approx 0.21$ 정도의 잔차가 불가피하다. $\ell^\infty$는 이 잔차에 의해 지배되어, 잘 형성된 형성에서도 낮은 Bind 값을 준다. $\ell^2$는 이를 전체 지점의 평균 잔차로 완화한다.

\textbf{증명된 부등식}: Cauchy--Schwarz 부등식으로부터,
\begin{equation}
    \mathsf{Bind} \geq 1 - \sqrt{\frac{\mathcal{E}_{\mathrm{cl}}}{n}}
    \label{eq:bind-energy-bound}
\end{equation}
즉, 작은 닫힘 에너지는 높은 Bind를 함의한다.


\section{Sep: 외부와의 구조적 대비}
\label{sec:sep}

\subsection{직관}

``이 형성은 그 주변과 구분되는가?''

Sep(분리)는 형성의 내부가 외부와 구조적으로 얼마나 비대칭인지를 측정한다. 높은 결합(Bind)만으로는 충분하지 않다---자기-지지적이지만 주변과 구분되지 않는 장(예: 균일한 장 $u \equiv c$)은 형성이 아니다. 형성은 자신의 외부와 \emph{다르게} 조직되어야 한다.

비유하자면: 세포를 생각하자. 세포의 내부(세포질)는 자기-지지적이지만, 그것만으로 세포가 ``하나의 것''이 되는 것은 아니다. 세포가 하나의 개체로 존재하려면 외부(세포외액)와 구조적으로 \emph{구분}되어야 한다. Sep는 이 구분의 정도를 측정한다.


\subsection{수학적 정의}

구별 연산자 $\mathbf{D}_t(x; 1-u_t) \in [0,1]$는 지점 $x$가 외부 장 $1-u_t$에 대해 가지는 구조적 비대칭을 측정한다. 이를 이용하여:

\begin{equation}
    \boxed{\mathsf{Sep}_t(u_t) = \frac{\sum_{x \in X_t} u_t(x)\, \mathbf{D}_t(x;\, 1-u_t)}{\sum_{x \in X_t} u_t(x)}}
    \label{eq:sep-def}
\end{equation}

이것은 구별의 $u$-가중 평균이다. 응집 값 $u_t(x)$가 자연스러운 중요도 가중치 역할을 한다: 높은 응집을 가진 지점이 분리 점수에 더 많이 기여한다.

\textbf{해석}:
\begin{itemize}
    \item $\mathsf{Sep} = 1$: 모든 응집 지점이 외부와 완전히 비대칭
    \item $\mathsf{Sep} \approx 0.5$: 무작위 또는 구조 없는 장
    \item $\mathsf{Sep} \approx 0$: 응집 지점이 외부와 구별되지 않음
\end{itemize}

\textbf{증명된 등식}: 놀랍게도 Sep와 분리 에너지 사이에 정확한 등식이 성립한다:
\begin{equation}
    \mathsf{Sep} = 1 - \frac{\mathcal{E}_{\mathrm{sep}}}{m}
    \label{eq:sep-energy-identity}
\end{equation}
여기서 $m = \sum u_t(x)$는 총 응집 부피이다. 이 등식은 진단 요소와 에너지 항 사이의 직접적 다리(bridge)를 제공한다.

% [Figure placeholder: Sep visualization]
\begin{figure}[t]
\centering
% \includegraphics[width=0.8\textwidth]{figures/ch02_sep_visualization.pdf}
\fbox{\parbox{0.8\textwidth}{\centering\vspace{3.5cm}
\textbf{[그림 2.2]} Sep 시각화.\\
좌: 원본 장 $u$ (5$\times$5 격자의 bump). 중앙: 구별 $\mathbf{D}(x; 1-u)$ (내부에서 높음).\\
우: $u$-가중 구별 $u(x) \cdot \mathbf{D}(x; 1-u)$ (Sep 계산의 분자).
\vspace{0.5cm}}}
\caption{Sep 계산 과정. 구별 연산자(중앙)는 내부 지점에서 높은 값을, 외부 지점에서 낮은 값을 가진다. 응집 가중(우측)은 형성의 지지 영역에 자연스럽게 집중한다.}
\label{fig:sep-visualization}
\end{figure}

\textbf{왜 $u$-가중인가?} 이전 버전에서는 공동-소속(co-belonging) 연산자 $\mathbf{C}_t$를 가중치로 사용하는 형태가 고려되었다. 그러나 $\mathbf{C}_t(x,x)$는 외부 지점에도 상당한 가중치를 부여하여, $\mathbf{C}_t$-가중 평균은 형성의 질에 무관하게 항상 $\approx 0.5$를 준다는 것이 발견되었다. $u$-가중은 형성의 자체 지지 영역에 자연스럽게 집중하여 의미 있는 진단값을 제공한다.


\section{Inside: 형태학적 분절}
\label{sec:inside}

\subsection{직관}

``이 형성은 핵심부, 경계, 외부가 명확한 구조를 가지는가?''

Inside(내부성)는 형성이 단순히 존재하는 것을 넘어, 명확한 형태학적 구조---뚜렷한 핵심부(core), 일관된 경계 대역(boundary band), 그리고 깨끗한 외부(exterior)---를 가지는지를 평가한다.

이 요소를 이해하기 위해, 지형학적 비유를 사용하자. 응집 장 $u_t$를 산의 높이로 생각하면:
\begin{itemize}
    \item 잘 형성된 형성 = 뚜렷한 봉우리가 하나 있는 산: 정상(핵심부)에서 기슭(외부)까지 매끄럽게 하강
    \item 나쁜 형성 = 여러 개의 비슷한 봉우리가 있는 산맥: 단일 지배적 구조 없음
    \item 균일한 장 = 평원: 높이가 일정하여 형태학적 구조 없음
\end{itemize}


\subsection{수학적 정의}

Inside는 $H_0$ 지속적 상동성(persistent homology)을 통해 정의된다. 핵심 아이디어는 상위-수준 집합 여과(superlevel-set filtration)이다: 임계값 $\theta$를 1에서 0으로 내리면서, $u_t(x) \geq \theta$인 지점들의 연결 성분을 추적한다.

이 과정에서 두 가지 핵심 양이 계산된다:
\begin{itemize}
    \item $\ell_{\max}$: 가장 긴 막대(bar)의 길이 = 지배적 연결 성분의 출생-사망 구간
    \item $\ell_{\mathrm{second}}$: 두 번째로 긴 막대의 길이 = 이차적 구조의 규모
\end{itemize}

이로부터 형태학적 질 측도(morphological quality measure)를 정의한다:
\begin{equation}
    \boxed{\mathsf{Inside}_t(u_t) = \mathcal{Q}_{\mathrm{morph}}(u_t) = \frac{\ell_{\max} - c}{1 - c} \cdot \mathrm{Artic}(u_t)}
    \label{eq:inside-def}
\end{equation}

여기서:
\begin{itemize}
    \item $c = m/n = \sum u_t / n$: 부피 분율 (장의 평균값)
    \item $\mathrm{Artic}(u_t) = 1 - \ell_{\mathrm{second}} / \ell_{\max}$: 분절 비율
    \item $(\ell_{\max} - c)/(1-c)$: 기준선(baseline) 보정---균일 장 ($\ell_{\max} = c$)에서 0을 보장
\end{itemize}

\textbf{해석}:
\begin{itemize}
    \item $\mathsf{Inside} \approx 1$: 단일 지배적 구조, 긴 주 막대, 이차 구조 미미
    \item $\mathsf{Inside} \approx 0$: 균일하거나 ($\ell_{\max} \approx c$) 여러 경쟁 구조 ($\mathrm{Artic} \approx 0$)
\end{itemize}

이 정의는 두 가지 독립적 요구를 곱의 형태로 포착한다: 정규화된 지속성은 일차 형성이 기준선 위로 얼마나 강건한지를, 분절 비율은 단일 형성이 장을 얼마나 깨끗하게 지배하는지를 측정한다.

% [Figure placeholder: Inside / persistence diagram]
\begin{figure}[t]
\centering
% \includegraphics[width=0.85\textwidth]{figures/ch02_inside_persistence.pdf}
\fbox{\parbox{0.85\textwidth}{\centering\vspace{4cm}
\textbf{[그림 2.3]} Inside 계산 과정.\\
좌: 응집 장 $u$의 상위-수준 집합 여과 (임계값 $\theta = 0.9, 0.7, 0.5, 0.3, 0.1$).\\
우: $H_0$ 지속성 다이어그램. 가장 긴 막대 $\ell_{\max}$와 두 번째 막대 $\ell_{\mathrm{second}}$가 표시됨.\\
$\mathrm{Artic} = 1 - \ell_{\mathrm{second}}/\ell_{\max}$.
\vspace{0.5cm}}}
\caption{$H_0$ 지속적 상동성을 이용한 Inside 계산. 잘 형성된 bump(상단)는 하나의 긴 막대가 지배적이며 높은 Inside 값을 가진다. 이중 bump(하단)는 두 개의 비슷한 막대를 가져 낮은 Artic를 준다.}
\label{fig:inside-persistence}
\end{figure}


\subsection{왜 지속적 상동성인가?}

Inside를 정의하는 데 $H_0$ 지속적 상동성을 사용하는 것이 왜 자연스러운가?

\begin{enumerate}
    \item \textbf{임계값 독립적}: 특정 임계값 $\theta$에 의존하지 않고, 모든 임계값에 걸친 위상적 구조를 포착한다.
    \item \textbf{구조 민감적}: 단순한 통계량(평균, 분산)으로는 구분할 수 없는 형태학적 차이(단일 bump vs.\ 이중 bump 등)를 포착한다.
    \item \textbf{이론적 일관성}: 부피 제약 하에서 균일 장에 대해 $\mathcal{Q}_{\mathrm{morph}} = 0$이 되는 공리(QM1)를 자연스럽게 만족한다.
\end{enumerate}


\section{Persist: 시간적 구조 계승}
\label{sec:persist}

\subsection{직관}

``이 형성은 이전 시각의 형성을 계승하는가?''

Persist(지속성)는 시간적 차원을 평가한다. 아무리 잘 결합되고, 잘 분리되고, 잘 분절된 형성이라도, 그것이 시간적 계보(temporal lineage) 없이 갑자기 나타난 것이라면 완전한 형성이 아니다. 진정한 형성은 이전 시각의 응집 조직을 구조적으로 계승한다.

비유하자면: 의식의 흐름에서 ``나''라는 인식이 매 순간 새로 생겨나는 것이 아니라, 이전 순간의 ``나''로부터 연속적으로 이어지는 것처럼, 형성의 정체성도 시간적 계승에 기반한다. Persist는 이 계승의 정도를 측정한다.


\subsection{수학적 정의: 핵심 중첩 근사}

SCC는 Persist에 대해 두 가지 구현을 제공한다. 여기서는 더 직관적인 핵심 중첩(core-overlap) 근사를 소개한다:

\begin{equation}
    \boxed{\mathsf{Persist}(u_{\mathrm{prev}}, u_{\mathrm{curr}}) = \frac{\sum_i \min(u_{\mathrm{curr},i},\, u_{\mathrm{prev},i})}{\max\!\big(\sum_i u_{\mathrm{curr},i},\, \sum_i u_{\mathrm{prev},i}\big)}}
    \label{eq:persist-overlap}
\end{equation}

이것은 두 연속 장 사이의 부드러운 중첩 측도이다:
\begin{itemize}
    \item \textbf{분자}: 지점별 최솟값의 합 = 두 장이 공유하는 응집의 총량
    \item \textbf{분모}: 두 장의 부피 중 큰 것 = 정규화 기준
\end{itemize}

\textbf{해석}:
\begin{itemize}
    \item $\mathsf{Persist} = 1$: 두 장이 동일 ($u_{\mathrm{prev}} = u_{\mathrm{curr}}$)---완벽한 지속
    \item $\mathsf{Persist} = 0$: 두 장이 서로소 지지(disjoint support)---지속 실패
    \item $0 < \mathsf{Persist} < 1$: 부분적 계승---형성이 이동하거나 변형됨
\end{itemize}

\textbf{정적 최적화}: 단일 시각만 고려하는 경우 ($u_{\mathrm{prev}}$ 없음), $\mathsf{Persist} = 1$로 설정한다.


\subsection{전달 기반 Persist}

보다 정교한 구현은 시간적 전달 핵 $\mathbf{M}_{t \to s}$를 이용한다. 핵심 아이디어는: 단순한 공간적 중첩이 아니라, 응집적 핵심의 \emph{구조적 계승}을 평가하는 것이다. 이 구현은 자기-참조적 비용 함수를 가진 엔트로피 정규화 부분 최적 전달(entropy-regularized partial optimal transport)에 기반한다. 자세한 내용은 시간적 구조를 다루는 Part~III에서 전개한다.


\section{예제: 5$\times$5 격자의 Gaussian Bump}
\label{sec:example-gaussian}

이론적 정의만으로는 직관이 불분명할 수 있다. 구체적 예제에서 각 진단 요소가 어떤 값을 가지는지 계산해보자.


\subsection{설정}

$5 \times 5$ 격자 그래프 $X = \{(i,j) : 0 \leq i,j \leq 4\}$를 고려한다. 총 $n = 25$개 지점이 있으며, 4-연결(상하좌우) 인접 구조를 가진다. 이 위에 중심 $(2,2)$에 놓인 Gaussian bump 형태의 응집 장을 정의한다:
\begin{equation}
    \tilde{u}(i,j) = \exp\!\left(-\frac{(i-2)^2 + (j-2)^2}{2 \cdot 1.0^2}\right)
    \label{eq:gaussian-bump-raw}
\end{equation}
이를 부피 제약 $\Sigma_m$에 사영하면 실제 응집 장 $u$를 얻는다.

% [Figure placeholder: Gaussian bump on 5x5 grid]
\begin{figure}[t]
\centering
% \includegraphics[width=0.6\textwidth]{figures/ch02_gaussian_5x5.pdf}
\fbox{\parbox{0.6\textwidth}{\centering\vspace{4cm}
\textbf{[그림 2.4]} $5 \times 5$ 격자 위의 Gaussian bump 응집 장.\\
색상: 응집 강도 (빨강=높음, 파랑=낮음).\\
중심 $(2,2)$에서 최대, 모서리에서 최소.
\vspace{0.5cm}}}
\caption{$5 \times 5$ 격자 위의 Gaussian bump 응집 장. 중심 지점이 가장 높은 응집을 가지며, 외곽으로 갈수록 점진적으로 감소한다. 이는 잘 구조화된 단일 형성의 전형적 형태이다.}
\label{fig:gaussian-5x5}
\end{figure}


\subsection{진단 벡터 계산}

Python \texttt{scc} 패키지를 이용한 계산 과정을 개략적으로 보여준다:

\begin{verbatim}
from scc import GraphState, ParameterRegistry, find_formation

# 5x5 격자 생성
graph = GraphState.grid_2d(5, 5)
params = ParameterRegistry()

# 형성 탐색 (에너지 최소화)
result = find_formation(graph, params)

# 진단 벡터 출력
print(result.diagnostics)
# DiagnosticVector(Bind=0.97, Sep=0.85, Inside=0.91, Persist=1.00)
\end{verbatim}

각 값의 의미를 해석해보자:

\begin{description}
    \item[$\mathsf{Bind} = 0.97$:] 장이 닫힘의 고정점에 매우 가깝다. 갱신해도 거의 변하지 않는다. 관계적 자기-지지가 거의 완벽하다.

    \item[$\mathsf{Sep} = 0.85$:] 내부 지점들이 외부와 상당히 구분된다. 완벽하지 않은 이유는 경계 대역의 지점들이 중간 정도의 구별 값을 가지기 때문이다.

    \item[$\mathsf{Inside} = 0.91$:] $H_0$ 지속성 분석에서 하나의 지배적 연결 성분이 확인된다. 높은 $\ell_{\max}$와 높은 $\mathrm{Artic}$가 곱해져 높은 값을 준다.

    \item[$\mathsf{Persist} = 1.00$:] 정적(단일 시각) 최적화이므로, 이전 장이 없어 $\mathsf{Persist} = 1.0$으로 설정된다.
\end{description}

\begin{table}[t]
\centering
\caption{$5 \times 5$ 격자 Gaussian bump의 진단 벡터 (대표적 수치)}
\label{tab:gaussian-diagnostics}
\begin{tabular}{lcc}
\hline
\textbf{요소} & \textbf{값} & \textbf{해석} \\
\hline
$\mathsf{Bind}$ & $\sim 0.97$ & 거의 완벽한 자기-지지 \\
$\mathsf{Sep}$ & $\sim 0.85$ & 양호한 외부 구분 (경계 대역 영향) \\
$\mathsf{Inside}$ & $\sim 0.91$ & 강한 단일 형태학적 구조 \\
$\mathsf{Persist}$ & $1.00$ & 정적 최적화 (기본값) \\
\hline
$\min(\mathbf{d})$ & $\sim 0.85$ & 형성의 최약점: 분리 \\
$\mathrm{mean}(\mathbf{d})$ & $\sim 0.93$ & 전반적으로 우수한 형성 \\
\hline
\end{tabular}
\end{table}


\subsection{비교: 균일 장과 무작위 장}

진단 벡터의 의미를 더 명확히 하기 위해, Gaussian bump 외의 두 가지 ``나쁜'' 장과 비교해보자.

\textbf{균일 장} $u \equiv c$ (모든 지점에서 동일한 값):
\begin{itemize}
    \item $\mathsf{Bind} \approx 1$: 균일 장은 닫힘의 고정점이다 (닫힘은 이미 균일한 장을 변화시키지 않음)
    \item $\mathsf{Sep} \approx 0.5$: 모든 지점이 외부와 대칭적---구분 없음
    \item $\mathsf{Inside} = 0$: $\ell_{\max} = c$이므로 정규화 후 0
    \item 결론: 높은 Bind에도 불구하고, Sep와 Inside가 낮아 형성이 아님
\end{itemize}

\textbf{무작위 장} (독립 균등 분포에서 추출 후 $\Sigma_m$에 사영):
\begin{itemize}
    \item $\mathsf{Bind}$: 낮음 (무작위 값은 관계적 자기-지지와 무관)
    \item $\mathsf{Sep}$: 중간 (임의적 구별)
    \item $\mathsf{Inside}$: 낮음 (여러 작은 연결 성분, 지배적 구조 없음)
    \item 결론: 모든 요소가 낮아 형성이 아님
\end{itemize}

이러한 비교는 진단 벡터의 네 요소가 \emph{독립적}으로 실패할 수 있음을 보여준다. 이것이 바로 단일 숫자가 아닌 벡터가 필요한 이유이다.


\section{종합: Proto-cohesion이란 무엇인가?}
\label{sec:synthesis}

이 장의 내용을 종합하자.

\textbf{Proto-cohesion}(원초-응집)은 응집 장이 진정한 응집적 형성으로서 가지는 질(quality)의 연속적 평가이다. 이것은 불리언 판정(``형성인가 아닌가?'')이 아니라, 네 개의 독립적 차원에서의 연속적 평가이다.

\begin{equation}
    \mathbf{d} = (\mathsf{Bind},\ \mathsf{Sep},\ \mathsf{Inside},\ \mathsf{Persist}) \in [0,1]^4
\end{equation}

이 벡터는 형성의 상태를 $[0,1]^4$의 한 점으로 표현한다. 이 4차원 공간에서:
\begin{itemize}
    \item $(1, 1, 1, 1)$: 이상적 객체---모든 구조적 요구가 완벽히 충족
    \item $(0, 0, 0, 0)$: 전적 비-형성---어떤 구조적 요구도 충족되지 않음
    \item 중간 지점: 부분적 형성---일부 요구는 충족되고 일부는 아닌 전-객체적 상태
\end{itemize}

전통적 접근은 $(1,1,1,1)$만을 ``물건''으로 인정하고 나머지를 무시한다. SCC는 전체 $[0,1]^4$ 공간을 이론의 대상으로 삼는다. 이를 통해 전-객체적 상태---형성되고 있지만 아직 완전한 객체는 아닌 상태---를 체계적으로 기술할 수 있다.

\begin{center}
\fbox{\parbox{0.85\textwidth}{\centering
\textbf{핵심 메시지}\\[6pt]
Proto-cohesion은 ``물건인가 아닌가''의 이분법을 대체한다.\\
그 자리에 ``어떤 측면에서 얼마나 응집적인가''의 연속적 진단을 놓는다.\\
이것이 SCC가 전-객체적 인식을 수학적으로 포착하는 방식이다.
}}
\end{center}

\bigskip

다음 장에서는 이 진단 벡터의 각 요소가 어떤 에너지 항과 대응하는지, 그리고 에너지 최소화가 어떻게 좋은 형성을 낳는지를 수학적으로 전개한다.
